\subsection{Implementasi Token Otorisasi Berbasis Macaroons}
\label{subsec:implementasi-token-macaroons}

\subsubsection{Komponen Implementasi Macaroons}
\label{subsubsec:komponen-implementasi-macaroons}

Seperti yang sudah dijelaskan pada subbab \ref{subsec:implementasi-arsitektur-sistem}, terdapat dua \textit{crate} yang berkaitan dengan implementasi Macaroons pada implementasi Tugas Akhir ini, yaitu \texttt{macaroon} dan \texttt{decmed-macaroon-auth}.

\textit{Crate} \texttt{macaroon} merupakan hasil penyesuaian dari implementasi Macaroons berbasis Rust yang tersedia secara \textit{open source} pada URL \url{https://github.com/macaroon-rs/macaroon}. Penyesuaian dilakukan karena \textit{crate} asli memerlukan beberapa penyesuaian agar lebih sesuai dengan kebutuhan implementasi sistem DecMed, terutama dari sisi kompatibilitas \textit{toolchain} dan pengelolaan \textit{dependency}. Perubahan utama pada \textit{crate} \texttt{macaroon} adalah penggantian \textit{dependency} kriptografi berbasis \texttt{sodiumoxide} dengan \textit{crate} pure Rust, yaitu \texttt{hmac}, \texttt{sha2}, \texttt{chacha20poly1305}, dan \texttt{rand}.

Setelah perubahan tersebut dilakukan, \textit{unit testing} pada \textit{crate} dijalankan untuk memastikan bahwa fungsi utama Macaroons tetap berjalan dengan benar. Pengujian tersebut mencakup pembuatan Macaroon, penambahan \textit{caveat}, serialisasi, deserialisasi, serta verifikasi token. Dengan demikian, \textit{crate} \texttt{macaroon} digunakan sebagai fondasi implementasi Macaroons pada sistem DecMed.

\texttt{decmed-macaroon-auth} merupakan \textit{wrapper} yang dibangun untuk mempermudah penggunaan \textit{crate} \texttt{macaroon} pada implementasi sistem DecMed. \textit{Wrapper} ini menyediakan fungsi-fungsi utilitas untuk kebutuhan otorisasi DecMed, seperti proses penerbitan token Macaroon, melakukan atenuasi dengan penambahan \textit{caveat} baru, definisi struktur \textit{caveats} Macaroon, verifikasi Macaroon, dan fungsi pendukung lainnya. Rincian modul dan implementasi yang diimplementasikan pada \textit{crate} \texttt{decmed-macaroon-auth} ditunjukkan pada Lampiran \ref{lampiran:implementasi-crate-decmed-macaroon-auth}.

\subsubsection{Struktur \textit{Caveats} Macaroons}
\label{subsubsec:struktur-caveats-macaroons}

Struktur \textit{caveats} pada token Macaroon pada implementasi ini ditunjukkan pada Tabel \ref{tab:implementasi-caveats-macaroons}. \textit{caveats} terdefinisi pada \textit{crate} \texttt{decmed-macaroon-auth}, sehingga \textit{hospital client} membentuk token dengan menambahkan \textit{caveats} sesuai dengan daftar \textit{caveats} yang terdefinisi pada \textit{crate}, dan PRE akan menolak \textit{request} jika terdapat \textit{caveats} yang tidak termasuk ke dalam daftar \textit{caveats}.

\begin{table}[!ht]
	\centering
	\caption{\textit{Caveats} yang Diimplementasikan pada Token Macaroon}
	\label{tab:implementasi-caveats-macaroons}
	\small
	\begin{tabular}{|p{0.34\textwidth}|p{0.56\textwidth}|}
		\hline
		\textbf{\textit{Caveat}}        & \textbf{Tipe Data}
		\\ \hline
		\texttt{patient\_address}       & \texttt{String} (alamat IOTA format heksadesimal)
		\\ \hline
		\texttt{related\_rme\_id}       & \texttt{String} (ID episode RME)
		\\ \hline
		\texttt{root\_subject}          & \texttt{String} (alamat IOTA format heksadesimal)
		\\ \hline
		\texttt{delegated\_by}          & \texttt{String} (alamat IOTA format heksadesimal)
		\\ \hline
		\texttt{delegated\_to}          & \texttt{String} (alamat IOTA format heksadesimal)
		\\ \hline
		\texttt{read\_dataset\_in}      & \texttt{HashSet<DatasetCategory>}
		\\ \hline
		\texttt{write\_dataset\_in}     & \texttt{HashSet<DatasetCategory>}
		\\ \hline
		\texttt{read\_function\_in}     & \texttt{HashSet<FunctionCategory>}
		\\ \hline
		\texttt{write\_function\_in}    & \texttt{HashSet<FunctionCategory>}
		\\ \hline
		\texttt{expires\_before}        & \texttt{DateTime<Utc>} (\texttt{YYYY-MM-DD HH:MM:SS})
		\\ \hline
		\texttt{max\_delegation\_depth} & \texttt{u32}
		\\ \hline
		\texttt{hospital\_cid}          & \texttt{String} (CID IPFS)
		\\ \hline
		\texttt{parent\_token\_hash}    & \texttt{String} (hash SHA-256 dari token induk)
		\\ \hline
		\texttt{purpose}                & \texttt{String} (nilai: \texttt{Read} atau \texttt{Update})
		\\ \hline
	\end{tabular}
\end{table}

\subsubsection{Penerbitan Token Macaroon}
\label{subsubsec:penerbitan-awal-token-macaroons}

Token otorisasi diimplementasikan menggunakan Macaroons sesuai rancangan pada Subbab \ref{subsec:rancangan-token-macaroons}. Pada implementasi ini, penerbitan token dilakukan oleh \textit{Server} PRE setelah permintaan akses dari pasien dinyatakan valid.

Pada proses penerbitan akses awal, \textit{Server} PRE menerbitkan dua jenis token untuk penerima akses, yaitu token \texttt{Read} dan token \texttt{Update}. Token \texttt{Read} memuat \textit{caveat} utama berupa \texttt{patient\_address}, \texttt{root\_subject}, \texttt{read\_dataset\_in}, \texttt{read\_function\_in}, \texttt{expires\_before}, \texttt{max\_delegation\_depth}, \texttt{hospital\_cid}, dan \texttt{purpose = Read}. Token \texttt{Update} memuat \textit{caveat} berupa \texttt{patient\_address}, \texttt{related\_rme\_id}, \texttt{root\_subject}, \texttt{write\_dataset\_in}, \texttt{write\_function\_in}, \texttt{expires\_before}, \texttt{max\_delegation\_depth}, \texttt{hospital\_cid}, dan \texttt{purpose = Update}.

Ketika melakukan penerbitan untuk petugas rekam medis, fungsi yang dipanggil adalah \texttt{issue\_admin\_personnel\_token} yang termasuk implementasi pada \textit{crate} \texttt{decmed-macaroon-auth}. Alur lengkap pemberian akses awal dari pasien kepada petugas rekam medis, termasuk pemindaian QR code, verifikasi identitas, pembentukan material PRE, penyimpanan ke Redis, dan pencatatan grant akses ke \textit{smart contract}, dijelaskan pada subbab \ref{subsec:implementasi-atenuasi-offline-delegasi}.

\subsubsection{Atenuasi Token Macaroon}
\label{subsubsec:atenuasi-token-macaroons}

Pada implementasi ini, pemegang token utama adalah para tenaga kesehatan, dengan role \texttt{AdministrativePersonnel} atau \texttt{MedicalPersonnel}. Pemegang token melakukan atenuasi token melalui \textit{hospital client} dengan melakukan \textit{request} atenuasi token ke \textit{server} PRE. Implementasi dari fungsi atenuasi ini terdapat pada \textit{crate} \texttt{decmed-macaroon-auth} dengan nama fungsi \texttt{attenuate\_macaroon}. Fungsi ini bekerja dengan tahapan melakukan deserialisasi token Macaroon, lalu menambahkan \textit{caveat} baru sesuai dengan \textit{caveats} yang ingin ditambahkan, kemudian melakukan serialisasi kembali terhadap token Macaroon yang telah diatenuasi sebelum di-\textit{return} sebagai \textit{response}. Perhitungan \textit{signature} baru dari Macaroon pada proses \textit{chained-HMAC} dilakukan pada setiap kali \textit{caveat} baru ditambahkan.

\subsubsection{Revokasi Token Macaroon}
\label{subsubsec:revokasi-token-macaroons}

Revokasi token pada implementasi ini sesuai dengan rancangan pada Subbab \ref{subsec:rancangan-penegakan-otoritas}. Implementasi revokasi akses melibatkan pencatatan di IOTA \textit{on-chain}, Redis PRE, dan \textit{server} PRE.

Setiap token turunan menyimpan keterkaitan terhadap token induknya melalui \textit{caveat} \texttt{parent\_token\_hash}. Informasi ini digunakan untuk mendukung pencabutan berantai, sehingga pencabutan terhadap token pada tingkat tertentu dapat berdampak pada token-token turunan yang bergantung kepadanya.

Pemeriksaan pencabutan dilakukan oleh \textit{Server} PRE sebelum akses terhadap data terenkripsi diberikan. Saat \textit{request} diterima oleh PRE, maka PRE akan memeriksa \textit{hash} dari token Macaroon yang digunakan, dan mengecek apakah ada di daftar \textit{revoked keys} yang tersimpan pada Redis PRE dengan pola \verb|revoked:token:<hash>|. Hal ini juga dilakukan dengan semua \textit{hash} token yang ada pada \textit{caveat} \texttt{parent\_token\_hash}. Jika ada salah satu \textit{hash} yang ditemukan, maka \textit{request} akan ditolak. Selain itu, proses ini juga akan menandai hak akses yang tercatat pada \textit{on-chain} sebagai \textit{revoked}.